\section{Solucion de problemas}

\subsection{No aparecen pacientes}

Revisar que HAPI este arriba y que Synthea haya cargado datos:

\begin{lstlisting}
docker ps
curl -sG http://localhost:8080/fhir/Patient \
  --data-urlencode '_summary=count' | jq '.total'
\end{lstlisting}

Si el conteo es cero, ejecutar nuevamente el job de poblacion:

\begin{lstlisting}
docker compose --profile local-hapi --profile local-translator --profile seed-data run --rm synthea-seed
\end{lstlisting}

\subsection{El health check aparece degradado}

Ejecutar:

\begin{lstlisting}
curl -s http://localhost:3000/api/v1/health/ready | jq
\end{lstlisting}

Si HAPI o el traductor aparecen abajo, revisar logs:

\begin{lstlisting}
docker compose logs api
docker compose logs hapi
docker compose logs cql-translator
\end{lstlisting}

\subsection{Una regla no genera card}

Revisar:

\begin{itemize}
  \item que la regla este guardada;
  \item que la validacion no tenga errores;
  \item que la expresion booleana exista, por ejemplo \texttt{Aplica};
  \item que el hook sea \texttt{patient-view};
  \item que la regla este publicada y activa si se quiere usar desde la ficha;
  \item que el paciente tenga los datos FHIR que la regla consulta;
  \item que el cambio se haya hecho en el mismo navegador/sandbox.
\end{itemize}

\subsection{Dos alumnos ven resultados distintos}

Eso es esperado. Cada navegador tiene un sandbox propio. Si dos alumnos toman
el mismo paciente base, sus reglas, cambios y cards no deberian mezclarse.

\subsection{Necesito empezar desde cero}

Desde el menu superior del sandbox, usar \texttt{Reiniciar sandbox}. Eso limpia
el trabajo del navegador actual sin afectar a otros alumnos.
